% Limitations of the study (Gate C, revised at the user's request): each limitation with its context and the data or methods
% that can resolve it, in the constructive form of the MaleCNS paper's Limitations; the sex of the model and the limits on
% generalizing it (the guide's SGBA rule).
Our results come from a model, so each statement about flies or patients is a prediction that awaits a test in the animal. We
built the model on the connectome of one male fly, and some cell types and connections differ between the
sexes\cellcite{berg2026,sturner2025}. Female connectomes of the brain and of the nerve cord came from separate
animals\cellcite{dorkenwald2024,azevedo2024}. A connectome that unites the brain and nerve cord of one female
fly\cellcite{bates2026} will allow the same training and tests in the other sex, once its proofreading and cell typing match
those of the male CNS. Each connectome is also a snapshot of one animal: cell types recur across individuals, with occasional
variation in neuron counts and connections\cellcite{schlegel2024}. Twins trained on several connectomes would separate what is
general in the mechanism from what is individual.

During the task, synapses change only by short-term depression at olfactory terminals, the model has no neuromodulation, and
its delays last under a second, shorter than the visual memories of several seconds in behaving
flies\cellcite{neuser2008,kuntz2017}. Adding short-term synaptic dynamics across the CNS and the dopamine-gated plasticity that writes associative memories in the mushroom body\cellcite{aso2016} would show whether the
intrinsic trace competes with synaptic memories over longer delays or complements them. The transplant between brain states
moves the full state, including its cue-independent part, so it separates where a trace was written from where it was read, not
the trace from the state that carries it. Transplants of single state variables can make that separation.

The molecular links rest on transcripts. Channel transcripts can track channel currents in identified
neurons\cellcite{schulz2006}, but disease can uncouple the two: in Aβ42 flies, \gene{Shal} transcripts rise\cellcite{park2025}
while the Kv4 channel is degraded\cellcite{ping2015}. The atlases map \val{M09} of the trace carriers, and most of the others lie
in the brain, where atlases that resolve lineages\cellcite{allen2026} or span the whole body\cellcite{li2022} can extend the
join. The disease operators are calibrated to published fly phenotypes rather than measured in the model's own cells.
Recordings from identified neurons of Aβ42 flies would calibrate them cell by cell, and the twin could then follow how such
shifts reshape memory as a disease progresses.
